\chapter{Testovanie softvéru}
\label{ch:kap2}
Cieľom tejto kapitoly je vymedziť pojem testovanie, jeho význam a dôsledky jeho zanedbania. Následne kapitola definuje termín Data-Driven testing (DDT), spôsob testovania, ktorý oddeľuje testovacie dáta od testovacej logiky a tým kladie zvýšené nároky na kvalitu a prípravu testovacích dát. Práve náročnosť prípravy dát je jeden z hlavných východísk tejto práce.

\section{Význam testovania}
Testovanie je možné chápať ako proces vyhodnocovania softvéru na základe pozorovania jeho behu [AmmannOffutt2016]. Jeho cieľom je odhaliť chyby a overiť, či softvér spĺňa špecifikované požiadavky a potreby používateľov~\cite{Myers2011}.

\section{Dôsledky nedôkladného testovania}
Každý softvér je v nejakej fáze vývoja chybný. Tieto chyby môžu mať za následok závažné straty, začínajúc od finančných strát až po straty na ľudských životoch. Existuje množstvo zaznamenaných prípadov kde nedôkladné testovanie viedlo až ku katastrofickým následkom. Jeden z týchto prípadov je explózia rakety  Ariane 5 (let 501), 37 sekúnd po jej štarte v roku 1996. Dôvodom bola neošetrená výnimka pri konverzii 64-bitového čísla s plávajúcou desatinnou čiarkou (float) na 16-bitové celé číslo so znamienkom (signed integer) v navigačnom systéme. Dôvodom bolo, že prevzatý modul nebol testovaný pre letové parametre novej rakety. Toto viedlo k strate odhadovanej na približne 370 miliónov dolárov a ročnému zdržaniu programu Ariane 5~\cite{lions1996}.

Ďalším príkladom je Therac-25, prístroj určený na rádioterapiu, ktorý bol používaný v USA a Kanade od roku 1983 až do jeho stiahnutia v roku 1987. Chybný softvér spôsobil šesť prípadov masívneho predávkovania radiáciou, z toho tri úmrtia. Podobné chyby boli nájdené aj v predošlom modeli Therac-20, ten mal ale hardvérové poistky, ktoré pri rovnakej softvérovej chybe zabránili ožiareniu. Therac-25 ich nahradil čisto softvérovou kontrolou~\cite{Leveson1993}.

Podobných prípadov je ale viac. Poukazujú na dôležitosť testovania, nevyhnuteľné pre bezpečnosť softvéru. Okrem katastrofických zlyhaní spôsobuje nekvalitný softvér aj priebežné ekonomické náklady~\cite{krasner2022}.


\section{Úrovne testovania}
Testovanie prebieha na rôznych úrovniach, pričom každá overuje iné aspekty systému~\cite{Ammann2016}.

\begin{description}[leftmargin=0pt]
    \item[Jednotkové testovanie (unit testing)] testovanie najmenších jednotiek kódu ako funkcie a metódy.
    \item[Testovanie modulov (module testing)] testovanie modulu (kolekcie jednotiek) v izolácii. Hľadajú sa chyby v komunikácii jednotiek v rámci modulu. 
    \item[Integračné testovanie (integration testing)] overenie rozhrania a komunikácie medzi modulmi. 
    \item[Testovanie systému (system testing)] kontroluje softvér ako celok. Hľadajú sa chyby v návrhu a špecifikácii.
    \item[Akceptačné testovanie (acceptance testing)] overenie, či softvér spĺňa požiadavky zákazníka a či spĺňa svoj primárny účel.
\end{description}

Táto práca sa dotýka predovšetkým vyšších úrovní, keďže generované testovacie dáta slúžia na plnenie databáz pri integračnom a systémovom testovaní aplikácií nad nimi.

\section{Automatizované testovanie}
Pri väčších projektoch sa manuálne testovanie stáva neefektívnym, nielen z časového, ale aj ekonomického hľadiska. Používa sa preto automatizované testovanie, pri ktorom testy vykonávajú špecializované nástroje a frameworky. Tento spôsob testovania sa stáva nevyhnutnou súčasťou vývoja. Automatizované testy sa spúšťajú opakovane a s rôznymi vstupmi, čím sa príprava testovacích dát stáva samostatnou, opakovanou potrebou. Jedným z prístupov testovania, ktorý oddeľuje testovacie dáta od testovacej logiky, je DDT, na ktorý sa zameriava táto práca.

\section{Data-Driven testing (DDT)}
DDT je prístup k testovaniu, ktorý oddeľuje testovacie dáta od testovacej logiky. Základnou myšlienkou je, že rovnaký testovací skript možno spustiť s rôznymi sadami vstupných dát, čím sa zvyšuje pokrytie testov bez nutnosti duplikovania kódu~\cite{Laukkanen2006}.

\begin{figure}[htbp]
    \centering
    \includegraphics[width=0.58\linewidth]{obrazky-figures/sadydat.pdf}
    \caption{jedna logika x N dát = N behov}
    \label{fig:dsadydat}
\end{figure}


\subsection{Hlavné charakteristiky} 
\begin{description}[leftmargin=0pt]
    \item[Oddelenie dát od logiky] testovacie dáta sú uložené v externých súboroch ako napríklad~Excel, CSV, YAML, nie v kóde testu.
    \item[Znovupoužiteľnosť] rovnaké testy možno spustiť s rôznymi sadami dát.
    \item[Čitateľnosť] testovacie dáta musia byť ľahko čitateľné aj osobou bez znalostí programovania~\cite{Laukkanen2006}.
    \item[Jednoduchá údržba] možné pridanie nového testovacieho prípadu bez zmeny kódu.
\end{description}


DDT prináša mnoho výhod, ako zvýšenie pokrytia testov a jednoduchá údržba testov, má však aj svoje nevýhody. Je potrebné testovacie dáta neustále udržiavať a príprava kvalitných testovacích dát môže byť náročná. Kvalita testovania je limitovaná kvalitou testovacích dát, čo je motivácia pre automatizované generovania testovacích dát, ktorému sa venuje táto práca.
